% ============================================================================
% SEC04.TEX - Section 8.4: HUD Design
% ============================================================================

\section{HUD Design}

\begin{learningoutcomes}
\item Merancang HUD yang efektif
\item Mengimplementasikan health bar, score display
\item Membuat HUD yang tidak mengganggu gameplay
\end{learningoutcomes}

\subsection{Elemen HUD Umum}

\begin{itemize}
\item Health bar
\item Score/Points display
\item Ammo counter
\item Minimap
\item Objective indicators
\end{itemize}

\subsection{Contoh Health Bar}

\begin{lstlisting}[language={[Sharp]C}]
using UnityEngine;
using UnityEngine.UI;

public class HealthBar : MonoBehaviour
{
    public Slider healthSlider;
    public Image fillImage;
    public Color fullHealthColor = Color.green;
    public Color lowHealthColor = Color.red;
    
    public void SetHealth(float currentHealth, float maxHealth)
    {
        healthSlider.value = currentHealth / maxHealth;
        
        // Change color based on health
        fillImage.color = Color.Lerp(lowHealthColor, fullHealthColor, healthSlider.value);
    }
}
\end{lstlisting}

\begin{catatan}
HUD harus informatif tapi tidak mengganggu. Gunakan corner atau edge screen untuk placement. Pertimbangkan opacity untuk tidak menghalangi view.
\end{catatan}
